% ── Figure: Copula Comparison ─────────────────────────────────────────────
\begin{figure}[tp]
    \centering
    \includegraphics[width=\textwidth]{Fig-Copula-Comparison.pdf}
    \Description{Two-panel figure comparing dependence models for a four-asset portfolio. Panel (a) shows a grouped bar chart of KS pass rates for SPY, NVDA, JNJ, and JPM under four dependence models: SIM, Gaussian copula, Student-t copula, and vine copula. All three copula models achieve pass rates above 95 percent for all tickers, while SIM falls below 30 percent for JPM. Panel (b) shows a correlation heatmap where each off-diagonal cell displays the observed and Student-t copula simulated correlation side by side, demonstrating close agreement.}
    \caption{Dependence model comparison for the four-asset portfolio (SPY, NVDA, JNJ, JPM).
    Panel~(a): KS pass rates by ticker and dependence model. All three copula specifications
    preserve per-asset marginal fidelity above 95\%, while the SIM degrades to 27.5\% for JPM
    due to low factor loading ($R^2 = 0.48$). The dashed line marks the 95\% threshold.
    Panel~(b): each off-diagonal cell shows observed / Student-t copula simulated pairwise
    correlations. The Student-t copula reproduces the cross-asset dependence structure with
    a pairwise MAE of 0.039; the largest discrepancy is SPY-JNJ (0.47 vs.\ 0.41).}
    \label{fig:copula_comparison}
\end{figure}

In this study we developed a hybrid hidden Markov model that combines a fine-grained quantile-based state partition ($N=100$) with a Poisson jump-duration mechanism to generate synthetic equity excess growth rate paths. We benchmarked it against seven alternative generators on SPY using a multi-dimensional quality evaluation (KS/AD pass rates, Wasserstein-1 and Hellinger distances, kurtosis matching, ACF-MAE, and quantile coverage) and extended it to a 424-asset universe via both a Single-Index Model factor decomposition and copula-based dependence specifications. The results showed that no single model dominated all quality dimensions: GARCH captured temporal structure but failed distributional tests, the GRU learned dynamics but suffered variance collapse, and the standard HMM matched distributions but could not reproduce volatility clustering. HMM-WJ was the only model to improve meaningfully on both dimensions, occupying the Pareto frontier between distributional fidelity and temporal structure. The paragraphs below interpret why the jump mechanism produces this advantage, discuss the practical implications for synthetic data generation and risk management, and identify the limitations that bound these conclusions.

The comparison between HMM-NJ and HMM-WJ isolated the contribution of the Poisson jump-duration mechanism. Without jumps, the Markov chain exited tail states within one to two steps on average, producing an ACF of $|G_t|$ that decayed too quickly. Introducing jumps forced the model to dwell in high-volatility regimes for empirically realistic durations \cite{ryden_stylized_1998, bulla_stylized_2006}. The HSMM baseline provided a direct test of the alternative semi-Markov approach: despite using the same Student-t emissions, the HSMM's coarse $K=8$ state partition caused empirical dwell times to average only 1.1--1.8 steps, collapsing to memoryless behavior (82\% KS vs.\ 98\% for HMM-WJ; ACF-MAE at the i.i.d.\ floor). This advantage was independent of the emission distribution choice (Table~\ref{tab:emission_hsmm_s9} in Online Appendix~S9). Our framework resolves this resolution-persistence tradeoff by maintaining fine state resolution ($N=100$) for distributional fidelity while enforcing tail-state persistence externally via the jump mechanism. This tradeoff shaped the broader model comparison: no single model dominated all quality dimensions simultaneously. HMM-WJ occupied the Pareto frontier between distributional fidelity and temporal structure: it reduced ACF-MAE by 28\% relative to the i.i.d.\ floor while maintaining pass rates above 91\% of the Bootstrap ceiling. GARCH achieved the best temporal fidelity but catastrophic distributional failure (5.5\% KS); the GRU learned temporal dynamics effectively but suffered variance collapse (0.6\% KS). The value of HMM-WJ lay in avoiding each alternative's severe failure mode. Crucially, the volatility-clustering strength is continuously tunable via $\epsilon$, and the grid search selected the value that jointly optimized temporal and distributional quality.

These tradeoffs carry practical implications. Generative models that reproduce the stylized facts of financial markets enable stress testing under market regimes that are statistically plausible but not historically observed, addressing a key limitation of scenario libraries derived solely from empirical records. The interpretable regime structure, where each hidden state corresponds to a quantile-defined segment of the excess growth rate distribution, facilitates communication between quantitative analysts and risk managers; states can be labeled (e.g., \textit{crash}, \textit{bear}, \textit{neutral}, \textit{bull}, \textit{rally}) and linked to economic narratives. The model's computational efficiency further supports Monte Carlo-based risk metrics such as Value-at-Risk and Conditional Value-at-Risk at scale across large asset universes \cite{glasserman_monte_2003}. This computational efficiency stems from the avoidance of the Baum-Welch EM algorithm, made possible by the quantile-based state partition that explicitly assigns each observation to a hidden state, eliminates convergence sensitivity to initialization and reduces computational cost. This property matters for large-scale pipelines. Fitting the model to 424 assets requires only empirical counting and a two-parameter grid search, making it feasible to regenerate synthetic scenarios on a daily or weekly cadence as new data arrive. For synthetic data applications where privacy is a concern, such as sharing simulated portfolios or stress-test scenarios with external parties, the framework offers implicit privacy protection: synthetic paths are generated from the learned Markov chain and emission distribution rather than slightly modified copies of historical records. While this generative structure reduces the risk of memorization, the framework does not provide formal differential privacy guarantees \cite{dwork_differential_2006}; quantifying membership inference risk and integrating calibrated noise mechanisms remain directions for future work \cite{stadler_synthetic_2022}. Beyond the model itself, the multi-metric evaluation approach used here, combining KS/AD pass rates with Wasserstein/Hellinger distances and ACF-MAE, can be applied to any time series generator where temporal dependence is critical to downstream applications \cite{stenger_jdiq_ts_synthesis_2024}.

Several limitations qualify these conclusions. First, the model assumes stationarity of the transition matrix and jump hyperparameters across the full in-sample window. The out-of-sample ACF analysis (panel~d) provided direct evidence of this limitation. The IS-calibrated $(\epsilon^*,\lambda^*)$ underestimated volatility clustering persistence in the 2025 test window, consistent with the elevated macro uncertainty of that period. Structural breaks in market microstructure, such as those accompanying regulatory changes or the COVID-19 shock, may not be adequately captured by a single static transition matrix. Relatedly, the quantile-based state partition is fixed at estimation time and does not adapt to evolving volatility regimes, potentially overweighting the central mass relative to the tails during prolonged market stress. Third, the Single-Index Model extension imposes a linear, single-factor decomposition whose median IS KS pass rate of 66.7\% across 424 assets (well below the 95\% level expected under a correct model) confirmed that it could not fully capture asset-specific tail behavior, skewness, or sector dynamics \cite{sharpe_simplified_1963, fama_french_1993}. The copula comparison demonstrated that replacing the SIM with a Student-t copula restored marginal KS pass rates above 95\% for all tested assets while also reproducing pairwise correlations to within 4 percentage points on average (MAE of 0.039 on the $[-1,1]$ correlation scale); this confirmed that the HMM marginal models were not the bottleneck and that the multi-asset quality gap was attributable to the dependence specification. The architecture is modular, and substituting a copula for the SIM does not require modifying the underlying Markov chain machinery; however, copula estimation scales as $O(d^2)$ in the number of assets, which may limit applicability to very large universes where the SIM's $O(d)$ scaling remains advantageous. Finally, the out-of-sample evaluation covers a single 249-day window (full calendar year 2025), which limits the statistical power of out-of-sample inference; a rolling evaluation would provide a more robust assessment of generalization performance.

A separate structural concern is that the jump mechanism produces a bimodal ensemble. At $\epsilon^{*}=10^{-4}$, approximately 76\% of simulated paths contain no jump events and therefore exhibit no volatility clustering, while the remaining 24\% exhibit pronounced clustering (Online Appendix~S7). For applications that draw a small number of paths (e.g., a single stress-test scenario), this means the majority of draws will lack temporal structure entirely. Users requiring volatility clustering in every path should either condition on jump-containing paths or increase $\epsilon$ at the cost of marginal fidelity. The optimal $\epsilon^{*}$ sits at the lower boundary of the search grid, and extending the grid below $10^{-4}$ could reveal that even rarer jumps are preferred by the objective; future work should explore whether alternative objective formulations that penalize the fraction of jump-free paths could address this limitation. On the data side, the excess growth rate calculation subtracts a constant risk-free rate $r_f$ derived from STRIPS yields, but the effective Federal Funds rate varied from 0.09\% to 5.33\% over the 2014--2024 sample period. This time-varying discrepancy introduces a slow-moving bias into the growth rate series that the model may interpret as regime structure rather than interest-rate dynamics. Additionally, the 2014--2024 training window includes the March 2020 COVID-19 crash, which likely dominates the empirical kurtosis estimate and the transition matrix's tail entries; whether the model's calibration generalizes to periods without such an extreme event remains untested. Finally, the two-sample KS and AD tests assume i.i.d.\ observations, but each simulated path exhibits temporal dependence by construction. Volatility clustering induces positive autocorrelation in the tail-indicator process $\mathbf{1}(G_t \le x)$, reducing the effective sample size below $T$ and potentially inflating pass rates slightly relative to what a block-bootstrap calibration would produce. The Wasserstein-1 and Hellinger distances carry no such sampling-distribution assumption and produce a consistent model ordering, corroborating the fidelity conclusions independently of this caveat.
